% Clase 18 --- El flujo por una superficie ABIERTA (§2.2).
% "Esto lo veo como una especie de red de mariposas, ¿se entiende? Donde esto
%  sería el borde... Es cualquier superficie bordeada por C."
% Los dos paneles tienen la misma curva C y el mismo campo: lo único que cambia
% es qué superficie se elige para calcular el flujo. Y da lo mismo.
\begin{center}
\begin{tikzpicture}[scale=1]

  % =============== (a) la superficie plana ===============
  \begin{scope}
    \foreach \xx in {-1.5,-0.75,0,0.75,1.5} {
      \draw[figvecaux, figblue] (\xx,-1.5) -- (\xx,1.7);
    }
    \node[figlblaux, figblue, anchor=south east] at (1.9,1.5) {$\vec B$};

    \fill[figamber!20] (0,0) ellipse (1.75 and 0.55);
    \draw[figline, line width=1.4pt] (0,0) ellipse (1.75 and 0.55);
    \draw[figvec, line width=1.1pt] (1.75,0.05) arc (0:22:1.75 and 0.55);
    \node[figlbl, figblue, anchor=west] at (1.85,0.28) {$C$};
    \draw[figvecres, line width=1.2pt] (0,0) -- (0,1.05);
    \node[figlbl, figred, anchor=west] at (0.08,0.75) {$\hat n$};
    \node[figlbl, figamber, anchor=north] at (0,-0.78) {$S$ plana};

    \node[figlbl, anchor=north, align=center] at (0,-1.75)
      {(a) la superficie más natural};
  \end{scope}

  % =============== (b) la "red de mariposas" ===============
  \begin{scope}[xshift=6.0cm]
    \foreach \xx in {-1.5,-0.75,0,0.75,1.5} {
      \draw[figvecaux, figblue] (\xx,-1.5) -- (\xx,1.7);
    }

    \fill[figamber!20] (-1.75,0) .. controls (-1.5,-1.5) and (1.5,-1.5) ..
      (1.75,0) arc (0:-180:1.75 and 0.55);
    \draw[figamber, line width=1.1pt] (-1.75,0) .. controls (-1.5,-1.5) and
      (1.5,-1.5) .. (1.75,0);
    \draw[figaux, figamber] (0,0) ellipse (1.75 and 0.55);
    \draw[figline, line width=1.4pt] (-1.75,0) arc (180:360:1.75 and 0.55);
    \draw[figline, line width=1.4pt, dashed] (1.75,0) arc (0:180:1.75 and 0.55);
    \draw[figvecres, line width=1.2pt] (0,0) -- (0,1.05);
    \node[figlbl, figred, anchor=west] at (0.08,0.75) {$\hat n$};
    \node[figlbl, figamber, anchor=north] at (0,-1.35) {$S'$ abombada};

    \node[figlbl, anchor=north, align=center] at (0,-1.75)
      {(b) otra superficie con el \textbf{mismo} borde:\\[-0.2em]
       el flujo es el mismo};
  \end{scope}

\end{tikzpicture}\\[0.5em]
{\small\itshape Acá hay una diferencia con Gauss que conviene tener presente:
en Gauss la superficie es \textbf{cerrada} y no tiene borde, mientras que en
Faraday es \textbf{abierta} y su borde es justamente la espira. Faraday lo
pensaba como una red de mariposas: el aro es $C$ y la tela es $S$, y la tela
puede estar más o menos abombada sin que eso cambie nada. El flujo
$\Phi_B = \iint_S \vec B\cdot\hat n\,dS$ resulta el mismo para cualquier
superficie apoyada en el mismo borde, y la razón es que las líneas de $\vec B$
son cerradas: todo lo que atraviesa una tela atraviesa la otra ---si no fuera
así, la ley sería ambigua y no serviría---. Lo que sí hay que fijar es la
orientación: elegido un sentido de recorrido para $C$, la mano derecha da
$\hat n$, y ese par de convenciones ligadas es lo que después permite escribir
la ley con signo. En el flujo entran los tres ingredientes que Faraday había
ensayado por separado: el módulo de $\vec B$, el área y el ángulo.}
\end{center}
